% !TeX program = xelatex
\documentclass[12pt,a4paper]{article}
\usepackage[margin=2cm]{geometry}
\usepackage{fontspec}
\usepackage{amsmath,amssymb}
\usepackage{enumitem}
\setmainfont{Times New Roman}
\setlength{\parindent}{0pt}
\setlength{\parskip}{6pt}
\begin{document}
\begin{center}
{\Large\bfseries ĐỀ THI TUYỂN SINH VÀO LỚP 10 THPT}\\
{\bfseries Thành phố Hà Nội, năm học 2010--2011}\\
{\itshape Phần Đại số}
\end{center}

\textbf{Câu 1.} Cho biểu thức
\[
P=\frac{\sqrt{x}}{\sqrt{x}+3}
+\frac{2\sqrt{x}}{\sqrt{x}-3}
-\frac{3x+9}{x-9},
\qquad x\ge0,\ x\ne9.
\]
\begin{enumerate}[label=\alph*)]
\item Rút gọn $P$.
\item Tìm $x$ để $P=\dfrac13$.
\item Tìm giá trị lớn nhất của $P$.
\end{enumerate}

\textbf{Câu 2.} Giải bài toán sau bằng cách lập phương trình:

Một mảnh đất hình chữ nhật có độ dài đường chéo là 13 m và chiều dài lớn hơn chiều rộng 7 m. Tính chiều dài và chiều rộng của mảnh đất đó.

\textbf{Câu 3.} Cho parabol $(P):y=-x^2$ và đường thẳng $(d):y=mx-1$.
\begin{enumerate}[label=\alph*)]
\item Chứng minh rằng với mọi $m$, đường thẳng $(d)$ luôn cắt $(P)$ tại hai điểm phân biệt.
\item Gọi $x_1,x_2$ là hoành độ các giao điểm của $(d)$ và $(P)$. Tìm $m$ để
\[
x_1^2x_2+x_2^2x_1-x_1x_2=3.
\]
\end{enumerate}

\textbf{Câu 5.} Giải phương trình
\[
x^2+4x+7=(x+4)\sqrt{x^2+7}.
\]
\end{document}
